% ============================================================
% ch32.tex —— 第 32 章 椭圆曲线：曲线上长出的群
% ============================================================
\chapter{椭圆曲线：曲线上长出的群}

终章。一条三次曲线即将同时拴住三件看似无关的宝物：费马大定理、椭圆的周长、你手机的支付密码。而撬动这一切的支点简单到一句话：\textbf{过曲线上两点的直线，必交第三点}。这份豪横的来源：把直线代入 $y^2 = (\text{三次})$，得到关于 $x$ 的\textbf{三次}方程——三个根，两个名额已被占，第三个无处可逃（贝祖定理 $3\times1 = 3$ 的现场执行，第 30 章）。

\section{主角登场：名不副实的曲线}

\begin{dingyi}[椭圆曲线（本书工作版）]
形如
\[
y^2 = x^3 + ax + b \qquad \left(\text{要求 } 4a^3 + 27b^2 \ne 0\right)
\]
的曲线。比如 $y^2 = x^3 + 1$、$y^2 = x^3 - x$。
\end{dingyi}

先澄清这桩数学圈最著名的名字冤案：\textbf{椭圆曲线不是椭圆}。椭圆是二次曲线（圆的拉伸近亲，第 29 章）；这里是\textbf{三次}曲线。名字的来历藏在微积分里：算\textbf{椭圆的周长}要用积分
\[
\int\sqrt{\frac{1 - k^2t^2}{1 - t^2}}\;dt ,
\]
（第 18 章的堆积机器原样上岗——只不过被积函数是根号套根号）。这类积分叫\textbf{椭圆积分}，没有初等原函数（刘维尔 1834 年证明——\textbf{椭圆没有周长公式}，这是"圆有 $\pi$ 而椭圆没有 $\pi$"的严格版）；数学家研究它们的"反函数"，反函数的自变量恰好住在这种三次曲线上——于是"椭圆积分""椭圆函数""椭圆曲线"三代同堂，名字一脉误传至今。\textbf{一条与椭圆无关的曲线，顶着椭圆的名字，替椭圆背了周长算不出的锅}。

那条技术条款 $4a^3 + 27b^2 \ne 0$ 是"三次式无重根"的判别式版（$x^3 + ax + b$ 有重根 $\iff 4a^3 + 27b^2 = 0$，可由第 29 章"重根 $\iff \Delta = 0$"的思路推出）。为什么要排除重根？第 31 章习题里 $y^2 = x^3$（重根在 $x=0$）在原点长出"\textbf{尖点}"，图像有折角、不光滑——光滑是群律的安全底线（尖点处切线不明，第 15 章"可导性"的曲面版）。\textbf{定义的条款是为了定理的寿命}——序章方法论的最后一次执行。

\begin{tikzpicture}[scale=0.85]
  \draw[->] (-1.8,0) -- (2.9,0) node[right] {$x$};
  \draw[->] (0,-3.2) -- (0,3.7) node[above] {$y$};
  \draw[thick,shenlan,domain=-1:2.08,smooth,samples=90] plot (\x,{sqrt(\x*\x*\x+1)});
  \draw[thick,shenlan,domain=-1:2.08,smooth,samples=90] plot (\x,{-sqrt(\x*\x*\x+1)});
  \draw[cheng!85!black] (-1,0) -- (2.08,3.12);
  \fill (0,1) circle (1.2pt) node[above left] {\scriptsize $A(0,1)$};
  \fill (2,3) circle (1.2pt) node[right] {\scriptsize $B(2,3)$};
  \fill (-1,0) circle (1.2pt) node[below left] {\scriptsize 第三点 $C(-1,0)$};
  \node at (1.5,-3.6) {\scriptsize $y^2 = x^3+1$：一个"鼻结"加两条臂；橙线为弦 $AB$};
\end{tikzpicture}

\section{弦切游戏：必然的第三点}

主角取 $y^2 = x^3 + 1$——干净的整数点一抓一把：$x = 0 \Rightarrow y = \pm1$；$x = 2 \Rightarrow y = \pm3$；$x = -1 \Rightarrow y = 0$。开玩。

\textbf{第一局（弦）}。取 $A = (0,1)$、$B = (2,3)$，连线斜率 $\tfrac{3-1}{2-0} = 1$，直线 $y = x+1$。代入曲线：
\[
(x+1)^2 = x^3 + 1 \;\Longrightarrow\; x^2 + 2x + 1 = x^3 + 1 \;\Longrightarrow\; x^3 - x^2 - 2x = 0 .
\]
分解：$x(x^2 - x - 2) = x(x-2)(x+1) = 0$——三个根：$x = 0$（$A$ 本人）、$x = 2$（$B$ 本人）、$x = -1$。\textbf{被迫多出的第三点}：$y = -1 + 1 = 0$，即 $C = (-1, 0)$！每一小步都只用了初中代数：代入、展开、因式分解。而"第三个根必然存在"不是运气——三次方程三个根（第 11 章代数基本定理），两个名额被 $A,B$ 占死，第三个无处可逃。

\textbf{第二局（切线）}。$A + A$ 的情形（"连线"退化为 $A$ 处的切线，$A$ 占两票）：切线斜率用隐函数求导——两边对 $x$ 求导（右边 $3x^2$，左边 $2y\cdot y'$，$y'$ 是 $y$ 对 $x$ 的导数，链式法则）：
\[
2y\,y' = 3x^2 \;\Longrightarrow\; y' = \frac{3x^2}{2y} .
\]
在 $A = (0,1)$ 处：$y' = \tfrac{0}{2} = 0$——切线水平：$y = 1$。代入：$1 = x^3 + 1$，即 $x^3 = 0$——\textbf{三重根全在 $x = 0$}！切线与曲线在 $A$ 处"三阶接触"（贴得比相切还紧一档）。第三点还是 $A$ 自己。

\textbf{第三局（竖线）}。过 $(0,1)$ 与 $(0,-1)$ 的"连线"是竖线 $x = 0$：代入 $y^2 = 1$，两点 $y = \pm1$——只有两票，第三票在哪？答：在\textbf{无穷远点 $O$}。把曲线补成齐次坐标（第 31 章）：$Y^2Z = X^3 + Z^3$，与 $Z = 0$（无穷远直线）联立得 $X^3 = 0$，即 $O = [0:1:0]$（竖直方向的无穷远点），重数一。\textbf{每条竖线都过 $O$}——竖线的第三票，永远记在 $O$ 名下。这个 $O$ 就是群律的"零"。

\section{从传送门到群：定义曲线上的加法}

传送门已通，现在装运算。规定（先照做，马上验证）：
\begin{jielun}
$A + B := $ 连线第三点 $C$ 沿 $x$ \textbf{轴的镜像}（$C = (x_c, y_c) \mapsto (x_c, -y_c)$）。
\end{jielun}
加镜像的理由藏在第三局里：竖线过 $P$ 与 $P' = (x,-y)$（镜像对）时第三点是 $O$——即"\textbf{$P$ 加 $P'$ 等于零}"，翻译过来：\textbf{$P$ 的加法逆元就是它的镜像}。镜像规则让"逆元"这条群公理自动免费。

立即试算。第一局的账：$A + B = $ 镜像$(C) = $ 镜像$(-1, 0) = (-1, 0)$（$C$ 在轴上，镜像是自己）。第二局：$A + A = $ 镜像$(A) = (0, -1)$。验算自洽性：$(0,-1)$ 是 $A$ 的镜像，故 $A + A + (0,-1) = (A+A) + (A \text{ 的逆}) = (0,-1)+(0,-1)$……直接算 $(0,-1)+(0,-1)$：切线在 $(0,-1)$ 处斜率 $\tfrac{3\cdot0}{2\cdot(-1)} = 0$，水平线 $y = -1$，代入 $1 = x^3+1$ 又是三重根 $x=0$，第三点 $(0,-1)$ 自己，镜像后 $= (0,1) = A$。于是 $(0,-1)+(0,-1) = A$，恰好与"$A+A = (0,-1)$"互为逆运算——\textbf{小群已经成形}。

\textbf{四项体检}（第 7 章的表，最后一次拉出来用）：
\begin{itemize}
\item \textbf{封闭} \checkmark：加完还在曲线上（第三点是曲线与直线的交点，镜像也在曲线上——$y$ 换 $-y$ 方程不变）；
\item \textbf{单位元} \checkmark：$O$（$P + O = P$：过 $P$ 与 $O$ 的"直线"是竖线，第三点是 $P$ 的镜像，再镜像回 $P$）；
\item \textbf{逆元} \checkmark：$-P = $ 镜像$(P)$（竖线第三点是 $O$）；
\item \textbf{交换律} \checkmark：连线不分先后，$A + B = B + A$；
\item \textbf{结合律}：\textbf{成立但深藏功名}——几何验证要动用第 30 章"九点定三次曲线"的自由度账（帕斯卡定理的近亲），是这门学科的著名计算；本书授予证书，不现场表演。
\end{itemize}
\textbf{一条几何曲线，肚子里长出了一个完整的群。}而且这次是"无限位的群"：有理点（坐标为分数的点）构成子群——\textbf{莫德尔定理}（1922）说这个子群由有限多个"发电机"生成（结构像 "$r$ 台独立发电机 $\oplus$ 一个有限小拧巴群"，发电机台数 $r$ 叫\textbf{秩}）。"数一条曲线上的有理点"从此成为数论的一门正经产业。

\begin{faming}
\textbf{土名}：弦切加法、第三点传送门 $\to$ \textbf{正式名}：椭圆曲线的\textbf{群律}（group law），1639 年被牛顿注意到雏形，十九世纪被彻底澄清。它是篇二（群）与篇六（方程的画）的合龙点——全书的两条主线在最后一条曲线上打了个结。
\end{faming}

\section{三件宝物}

\paragraph{宝物一：费马大定理的一角。}费马在《算术》页边写下的著名宣告（1637）："将一个立方数分成两个立方数之和……或任一高于二次的幂分成两个同次幂之和，都是不可能的。我发现了一个真正绝妙的证明，可惜页边太窄写不下。"其中 $n = 3$ 的情形（不存在正整数解的 $x^3 + y^3 = z^3$）等价于：曲线 $x^3 + y^3 = 1$ 上没有"合格的"有理点——而这条曲线换个坐标就是一条椭圆曲线（$x^3 + y^3$ 型与 $y^2 = x^3 + \text{常数}$ 型可互相变换，配方式 $Y^2 = X^3 - 432$）。\textbf{数论的地牢问题，变成了曲线上的点查户口}：欧拉证 $n=3$（1770，用类椭圆曲线的分解）、热尔曼推进奇素数情形、库默尔发明"理想数"横扫一大批……人类围攻 358 年，1994 年怀尔斯（Taylor 协助补上最后一洞）终结全局——路线正是：证明\textbf{谷山--志村猜想}（每条有理椭圆曲线都由模形式"生成"），而弗雷的机智在于：若费马有反例，能造出一条"不模"的椭圆曲线，矛盾。\textbf{椭圆曲线是三百年战役的主战场}。

\paragraph{宝物二：密码。}第 6 章的 RSA 靠"分解大数极难"守门。接班人 \textbf{ECC（椭圆曲线密码）}换了一道门：在（模大素数 $p$ 的）世界取曲线上一点 $P$ 与大数 $k$，算 $P + P + \cdots + P$（$k$ 次"弦切加法"，快速幂取模的原班技术，第 6 章）飞快；但从终点 $kP$ 反推 $k$（"椭圆曲线离散对数"），目前没有像样的算法。\textbf{同样的安全强度，ECC 的钥匙只有 RSA 的六分之一长}：256 位 ECC $\approx$ 3072 位 RSA。你手机里的支付通道、比特币的地址体系（secp256k1 曲线 $y^2 = x^3 + 7$）、大量 HTTPS 证书——多半是这条曲线在站岗。第 5 章的韩信、第 6 章的费马，到本章完成接力。

\paragraph{宝物三：BSD 猜想。}椭圆曲线的"秩"（发电机的个数，曲线有理点的身高）怎么算？Birch--Swinnerton-Dyer 猜想（1960 年代）给出了惊人的答案：\textbf{由该曲线对应的 $\zeta$ 型函数（$L$ 函数）在 $s = 1$ 处的性态读出}——零点的阶恰好等于秩。\textbf{又回到篇四的 $\zeta$ 家族}（第 24 章的直系后代）。它是千禧年七大难题之一，悬赏一百万美元，目前最好的成绩是 Coates--Wiles、Gross--Zagier、Kolyvagin 等人在特殊情形下的部分确认。从第 22 章的齿轮出发，到这里你已能看懂题面——这就是"数学之网"的成色。

\begin{dongshou}
\begin{itemize}
  \yisi 在 $y^2 = x^3 + 1$ 上完成第二局的延伸：算 $B + B$（$B = (2,3)$：切线斜率 $\tfrac{3\cdot4}{2\cdot3} = 2$，切线 $y = 2x - 1$；代入 $(2x-1)^2 = x^3 + 1$，即 $x^3 - 4x^2 + 4x = x(x-2)^2 = 0$——切点二重根 $x = 2$ \checkmark，第三点 $x = 0$、$y = -1$；镜像后 $B + B = (0, 1) = A$）。再算 $(0,-1) + (0,-1) = A$（正文已演），确认 $\{O,\ (0,1),\ (0,-1)\}$ 三点\textbf{自成一国}（封闭）——一个三元循环小群。那么 $(-1,0)$ 加进来会怎样？下一题见分晓。
  \yier 算 $A + (-1,0)$（$A = (0,1)$）：连线斜率 $\tfrac{0-1}{-1-0} = 1$，直线 $y = x+1$；联立得 $x(x-2)(x+1) = 0$，两根 $0, -1$ 已占，第三根 $x = 2$、$y = 3$——第三点 $(2,3)$，镜像后 $A + (-1,0) = (2, -3)$。\textbf{新点冒出来了}：三元小群不是全部。把 $(2, \pm3)$ 也纳入，验证六个点 $\{O,\ (0,\pm1),\ (-1,0),\ (2,\pm3)\}$ 两两相加不出这六点之外——这是完整的六元循环群（$y^2 = x^3+1$ 的全部有理点，莫德尔定理的"秩 $0$"实例）。
  \yier 验证 $y^2 = x^3 - x$（$a = -1, b = 0$）通过技术条款：$4a^3 + 27b^2 = -4 \ne 0$ \checkmark；并找出它的全部"整点"：$(0, 0), (\pm1, 0)$（$y^2 = x(x-1)(x+1)$，$y \ne 0$ 时需要三个相邻偶数之积为平方数，小范围试几组感受难度）。
\end{itemize}
\end{dongshou}

\begin{shaonao}
\textbf{椭圆周长的终审。}第 18 章的后门（牛顿--莱布尼茨）对圆周长完美收官：$\int_0^{2\pi}\sqrt{1 + \cos^2t}\,dt$ 型积分在圆退化为 $2\pi$。对椭圆呢？

(1) 椭圆 $\tfrac{x^2}{a^2} + \tfrac{y^2}{b^2} = 1$ 的周长积分（参数化 $x = a\cos t,\ y = b\sin t$，弧长微分第 25 章烧脑时刻的 $ds$）：
\[
L = \int_0^{2\pi}\sqrt{a^2\sin^2t + b^2\cos^2t}\;dt = 4a\int_0^{\pi/2}\sqrt{1 - e^2\sin^2 t}\;dt ,
\]
其中 $e = \sqrt{1 - b^2/a^2}$ 是离心率。

(2) $e = 0$（正圆）：积分退化为 $4a\cdot\tfrac\pi2 = 2\pi a$ \checkmark 圆周长。$e > 0$：根号套正弦，\textbf{无初等原函数}（刘维尔定理）——这就是"椭圆积分"的本尊。

(3) 拉马努金（1914）给出神级近似：$L \approx \pi\left[3(a+b) - \sqrt{(3a+b)(a+3b)}\right]$。用地球轨道近圆的情形（$a = 1, b = 0.999$）验证精度；再用极端椭圆（$a = 1, b = 0.1$）与数值积分对账，看拉马努金公式精确到几位。

\textbf{一条曲线的名字，一桩"没有公式"的公案，一场三百年接力}——本章的三个名字（椭圆积分、椭圆曲线、BSD）全部从这道积分发源。全书的最后一块拼图，恰好回到微积分的起点。
\end{shaonao}

\begin{chengshi}
结合律的证明（用"三次曲线过九点"的计数论证或代数直接展开）放行大学（Silverman--Tate 第一、二章，附录 D）；"$y^2 = x^3+1$ 的有理点群是六阶循环群"（正文习题的小群算不完——完整答案是六点：$\{O, (0,\pm1), (-1,0), (2,\pm3)\}$，恰好六个，循环），秩 $= 0$ 是它的"个性"；绝大多数曲线秩 $\ge 1$（有理点无穷多）。莫德尔定理、谷山--志村、BSD 各是一章大学教材，本章只负责把它们挂上网。ECC 的工程细节（曲线选择、侧信道攻击防护）是另一门学科，数学内核已如正文。
\end{chengshi}
